\documentclass[12pt]{article}
\usepackage[utf8]{inputenc}
\usepackage[T1]{fontenc}
\usepackage{hyperref}
\title{Samtao Language Resources for K-12 Curriculum Development}
\author{Compiled by Claude Sonnet 5.5 (claude-sonnet-5-5) for Andr\'as Kornai}
\date{}
\begin{document}
\maketitle

\section{Introduction}
Samtao (ISO 639-3 \texttt{stu}, Glottocode samt1238) is an Austroasiatic language spoken mainly in eastern Shan State, Myanmar. Population estimates in search results vary widely (about 11,000 in Myanmar and about 15,000 to 16,500 overall), well below the 68,000 in the task description. Web search found very little for this language.

\section{Primary Educational Resources}
No primers, textbooks or school programmes in Samtao were found.

\section{Linguistic and Reference Materials}
\begin{itemize}
\item Numeral data for Samtao from the Max Planck Institute for Evolutionary Anthropology: \url{https://lingweb.eva.mpg.de/channumerals/Samtao.htm}
\end{itemize}
No dictionary or grammar was found.

\section{Religious and Community Resources}
\begin{itemize}
\item Scripture Earth listing for Samtao (reports a Bible, YouVersion links and Global Recordings audio; I could not verify its current status): \url{https://www.scriptureearth.org/00eng.php?iso=stu}
\item Asia Harvest profile of the Samtao in Myanmar (lists no scripture or recordings): \url{https://operation.asiaharvest.org/MYANMAR-Samtao.pdf}
\item Joshua Project, Samtao in Myanmar: \url{https://www.joshuaproject.net/index.php/people_groups/14680/BM}
\end{itemize}

\section{Digital and Audio Resources}
Audio is reportedly listed on the Scripture Earth page above. I did not independently confirm it.

\section{Curriculum Development Context}
Most Samtao in Myanmar also speak Burmese and Shan. No evidence of Samtao-medium schooling was found. A curriculum would need new community-led orthography and materials, with the scripture material above as the only possible starting text.

\end{document}